For maximum performance, modern GPU programs require not only SIMT-style parallelism, but also software-managed concurrent compute and data movement.
Performance engineers must reason about subdividing work and data into the thread hierarchy (threads, warps, warpgroups, blocks, clusters) and safe usage of asynchronous tensor core and memcpy instructions that target each level of the memory hierarchy.
Unlike CPUs, which abstract away intra-thread instruction re-ordering, these asynchronous instructions expose explicit instruction reordering to the software.

To provide safety checking for synchronization and functional correctness while preserving low-level control,
we propose Exo-GPU, an imperative, array-based, user-schedulable language.
Constructs include sequential and hierarchical parallel loops, warp specialization, wrapped PTX assembly accelerator instructions (mma, wgmma, TMA), and split barriers, all explicitly inserted by the programmer.
For checking, \emph{deterministic} small step semantics dictate the expected sequential baseline program behavior.
The semantics yield numeric output values and a trace of load, store, and barrier events, each tagged with the threads and ``qualitative timeline'' (accelerator identifier) that execute the corresponding event in the lowered CUDA program.

Exo-GPU compilation breaks down into a \emph{rewrite} phase, where the program is optimized with user-selected rewrite rules that preserve numeric output values (but not traces),
and a \emph{sync-check} phase, where we reason, for concrete input sizes, that the barriers prevent any reordering of load/store events that could prevent the CUDA program from implementing the sequential baseline.
Together, the phases transitively ensure the CUDA program functionally matches the simpler pre-rewrite program.
\texttt{<results here>}.
